\documentclass[11pt]{article}
\usepackage[margin=1in]{geometry}
\usepackage[T1]{fontenc}
\usepackage[utf8]{inputenc}
\usepackage{lmodern,microtype,amsmath,amssymb,amsthm,mathtools,booktabs}
\usepackage[colorlinks=true,linkcolor=blue,citecolor=blue,urlcolor=blue]{hyperref}
\hypersetup{pdftitle={Gaussian conditioning below the mean-field scaling: a sharp fluctuation threshold},pdfauthor={Dmitry Vaintrob; Kaarel Hanni; OpenAI Codex}}
\setlength{\emergencystretch}{3em}
\newtheorem{theorem}{Theorem}[section]
\newtheorem{proposition}[theorem]{Proposition}
\newtheorem{lemma}[theorem]{Lemma}
\newtheorem{corollary}[theorem]{Corollary}
\theoremstyle{remark}
\newtheorem{remark}[theorem]{Remark}
\title{Gaussian conditioning below the mean-field scaling:\\a sharp fluctuation threshold}
\author{Dmitry Vaintrob \and Kaarel H\"anni \and OpenAI Codex}
\date{October 15, 2026}
\begin{document}
\maketitle
\begin{abstract}
For a Bayesian $\tanh$ network with a Gaussian readout, the hidden
feature kernel concentrates exponentially in width. When the scaling
exponent satisfies $\gamma<1$, this remains true after conditioning
on a fixed output margin at every Boolean input and on noiseless
training signs. For $0<\gamma<1$, a Gaussian moderate-deviation
principle determines the leading margin energy. We ask whether it
also determines the full conditioned posterior. For
$0<\gamma<1/2$, we prove relative approximation of every full-margin
orthant probability by a Gaussian with the deterministic finite-gain
kernel, and hence total-variation approximation of every Boolean
posterior. The threshold $1/2$ is sharp for a universal statement:
with only three input bits and one hidden layer, an explicit pair of
compatible labelings has identical Gaussian energy, but fourth-order
feature fluctuations change their posterior odds by
$\exp[(8586a^4/625)N^{2\gamma-1}+o(1+N^{2\gamma-1})]$.
At the original margin $a=1/2$, these corrections change a prediction
from less than $2/3$ confidence under the Gaussian comparison to more
than $2/3$ confidence under the actual network, already at
$\gamma=1/2$. Thus Gaussian leading rates hold throughout the
subcritical regime, while exact conditioned Gaussian behavior does not.
\end{abstract}

\section{Model, conditioning, and the question}

Let $\mathcal X=\{-1,1\}^d$. Append a constant coordinate,
$h^0(x)=(1,x)$, and use $L-1$ hidden layers of common width $N$:
\[
 h^s(x)=\tanh(W^s h^{s-1}(x)),\quad 1\le s<L,
 \qquad z_N(x)=W^Lh^{L-1}(x).
\]
There are no later biases and the readout is linear. All weights
are independent centered Gaussians, with variances
\[
 \operatorname{Var}(W^1_{ij})=N^{1-\gamma}/d,\qquad
 \operatorname{Var}(W^s_{ij})=N^{-\gamma}\ (1<s<L),\qquad
 \operatorname{Var}(W^L_i)=N^{-\gamma-1}.
\]
These are the conventions in the original notes\footnote{\url{https://kaarelh.github.io/ai/nnbayes-generalization.pdf}.},
continued to $\gamma<1$.
Let $Z_N=(z_N(x))_{x\in\mathcal X}$ and fix a margin $a>0$,
with $a=1/2$ the original choice.

The prior is first conditioned on $|Z_N(x)|>a$ at every input.
A training sample $S$ then conditions on exact signs. If
$\mathcal G_S$ denotes the compatible Boolean functions and
$O_g=\{z:g(x)z(x)>a\text{ for all }x\}$, the posterior is exactly
\[
 \pi_N(g\mid S)=
 \frac{\Pr\{Z_N\in O_g\}}
      {\sum_{h\in\mathcal G_S}\Pr\{Z_N\in O_h\}},
 \qquad g\in\mathcal G_S.
\]
We also write $A_g$ for the event $Z_N\in O_g$.
Thus the all-input margin normalization cancels, but its restriction
at unseen inputs remains in every orthant. For a target $f$, success
at an input $x$ means
\[
 \pi_N\{g:g(x)=f(x)\mid S\}>2/3.
\]
This is a sign posterior, not Gaussian regression with numeric
observations of the latent outputs.

Every result below fixes $d,L,a,\gamma$ before taking $N\to\infty$.
Training labels may be any fixed consistent collection. Constants
are not uniform as dimension grows. In particular, a comparison
proved here cannot silently be used in a joint dimension--width
limit or converted into a sample-complexity bound.

The exact Gaussian readout leaves just one source of non-Gaussian
output behavior: the randomness of its feature covariance. A rare
margin event can tilt that covariance by an amount tending to zero
while still changing posterior odds by an unbounded factor. This
is why convergence of the covariance and a Gaussian rate function
are insufficient by themselves to identify a conditioned posterior.

\section{Conditioned kernels and quantitative Gaussian comparison}
\label{sec:quantitative-comparison}

Fix the input dimension $d$, depth $L$, and margin $a>0$. Put
$M=2^d$, $b_N=N^{(1-\gamma)/2}$ and $t_N=N^\gamma$, and let $Q_N$
be the empirical covariance matrix of the last hidden layer. All norms
in this section are the maximum entry norm, unless another norm is
indicated. Constants may depend on $d,L,a$; none of the statements is
uniform in these parameters. The exact independent Gaussian readout gives
\[
 Z_N\mid Q_N\sim\mathcal N(0,Q_N/t_N).
\]
For a full Boolean labeling $g$, set
\[
 O_g=\{z:g(x)z(x)>a\text{ for every }x\}.
\]
For a collection $S$ of observed labels, let $\mathcal G_S$ be the
nonempty set of full labelings consistent with $S$, and let
$A_S=\bigcup_{g\in\mathcal G_S}O_g$. Thus conditioning on $A_S$
includes the margin requirement on every unobserved input.

\subsection{The finite-gain center and conditional concentration}

For a covariance matrix $R$, write
\[
 \Phi_b(R)_{xy}=\mathbb E[\tanh(bU_x)\tanh(bU_y)],
 \qquad U\sim\mathcal N(0,R).
\]
Define deterministic matrices
\[
 R_0(x,y)=\frac{1+x\cdot y}{d},\qquad
 K_N^1=\Phi_{b_N}(R_0),\qquad
 K_N^s=\Phi_{b_N}(K_N^{s-1}),\qquad
 \widehat K_N=K_N^{L-1}.
\]
The distinction between $\widehat K_N$ and $\mathbb E Q_N$ matters:
at depth greater than two the nonlinear covariance recursion does not
commute with expectation. The limiting saturated kernel is
\[
 K^{(0)}(x,y)=\frac{1+x\cdot y}{d+1},\qquad
 K^{(s+1)}=\frac2\pi\arcsin K^{(s)},\qquad K=K^{(L-1)},
\]
where the arcsine is applied entrywise. This matrix is strictly positive
definite. Indeed, the first-layer sign features associated to the
nonparallel augmented inputs $(1,x)$ are linearly independent: crossing
one of their distinct hyperplanes changes only its corresponding sign.
Their Gram matrix is therefore positive definite. Subsequent arcsine
maps preserve this property, by their positive power series and the
Schur product theorem.

\begin{lemma}[Concentration about the finite-gain kernel]
\label{lem:cmp-concentration}
For each fixed $\gamma<1$, there are positive constants $C,c,u_0$ and
$N_0$ such that, for $N\ge N_0$ and $0<u\le u_0$,
\begin{align}
 \Pr\{\|Q_N-\widehat K_N\|>u\}&\le C e^{-cNu^2},
 \label{eq:cmp-finite-concentration}\\
 \|\widehat K_N-K\|&\le C b_N^{-1}.
 \label{eq:cmp-finite-bias}
\end{align}
\end{lemma}

\begin{proof}
The diagonal variances of all limiting hidden kernels are one, and
their distinct-coordinate correlations have absolute value strictly
less than one. In fixed neighborhoods of these matrices, $\Phi_b$ has
a Lipschitz constant independent of $b$. To verify the uniformity, each
off-diagonal entry uses only a two-dimensional Gaussian marginal.
Differentiate its density with respect to its covariance parameters.
On a compact set of positive variances and correlations bounded away
from $\pm1$, the $L^1$ norms of these density derivatives are bounded.
The integrand $\tanh(bu)\tanh(bv)$ has absolute value at most one,
so the resulting derivative bound is independent of $b$. The same
one-dimensional argument handles diagonal entries. This argument does
not require the full input covariance $R_0$ to be invertible.

Given the preceding hidden layer, its $N$ successor rows are independent.
Each entry of their empirical covariance averages $N$ variables in
$[-1,1]$. Conditional Hoeffding bounds, a union bound over the $M^2$
entries, and the uniform Lipschitz estimate give
\eqref{eq:cmp-finite-concentration} by induction over the fixed number
of layers. At each stage one applies the previous-layer estimate with
a constant multiple of $u$; exits from the fixed neighborhoods have
probability at most $C e^{-cN}$ and are absorbed into the same bound.

For a centered Gaussian $U$ of variance bounded below,
\[
 \mathbb E|\tanh(bU)-\operatorname{sign}U|\le C/b.
\]
The Gaussian sign-correlation formula and the same Lipschitz induction
therefore give \eqref{eq:cmp-finite-bias}.
\end{proof}

\begin{proposition}[Conditioning preserves the frozen kernel]
\label{prop:cmp-conditional-freezing}
For every fixed $\gamma<1$,
\[
 Q_N\longrightarrow K
 \quad\text{in probability conditional on }A_S.
\]
For $0<\gamma<1$, there is $C_0$ such that
$\Pr(A_S)\ge e^{-C_0t_N}$ for all sufficiently large $N$.
Consequently, for fixed $\epsilon>0$,
\[
 \Pr\{\|Q_N-K\|>\epsilon\mid A_S\}
 \le C_\epsilon e^{-c_\epsilon N+C_0t_N}.
\]
More quantitatively, for a fixed sufficiently large constant $A$,
\[
 \Pr\left\{\|Q_N-K\|>
 C b_N^{-1}+A\sqrt{t_N/N}\,\middle|\,A_S\right\}
 \le C\exp[-(cA^2-C_0)t_N].
\]
These bounds may be chosen simultaneously for all observed-label
collections $S$ at the fixed input dimension.
\end{proposition}

\begin{proof}
Choose any $g\in\mathcal G_S$ and any fixed bounded box strictly inside
$O_g$. On a sufficiently small fixed neighborhood of $K$, all covariance
eigenvalues are bounded above and below. Integrating the density of
$\mathcal N(0,Q_N/t_N)$ over the box gives a lower bound
$c t_N^{M/2}e^{-Ct_N}$. That kernel neighborhood has probability tending
to one, proving the asserted denominator bound. Dividing the bounds of
Lemma~\ref{lem:cmp-concentration} by this denominator proves the claims
for positive $\gamma$. Here $t_N=o(N)$, and
$b_N^{-1}=\sqrt{t_N/N}$ tends to zero.

If $\gamma=0$, the normalized output converges to the nondegenerate
Gaussian $\mathcal N(0,K)$ and the conditioning event has a positive
limiting probability. If $\gamma<0$, its normalized margin is
$a\sqrt{t_N}\to0$; the limiting conditioning event requires only the
observed signs, again with positive Gaussian probability. The
unconditional kernel concentration proves the corresponding conditional
claim in both cases.
\end{proof}

\subsection{What holds at every positive subcritical exponent}

Set
\[
 I_K(z)=\frac12z^TK^{-1}z,\qquad
 E_K(g)=\min_{z\in\overline O_g}I_K(z),\qquad
 E_S=\min_{g\in\mathcal G_S}E_K(g).
\]
Each orthant has a unique minimizing vector $z_g$. Write
\[
 \mathcal H_S=\{g\in\mathcal G_S:E_K(g)=E_S\},\qquad
 V_S=\{z_g:g\in\mathcal H_S\}.
\]

\begin{proposition}[Conditional latent concentration and unique minima]
\label{prop:cmp-minimizers}
For every $0<\gamma<1$, the posterior concentrates on $V_S$. More
precisely, for every $\delta>0$ there is $\eta_\delta>0$ such that
\[
 \limsup_{N\to\infty}t_N^{-1}\log
 \Pr\{\operatorname{dist}(Z_N,V_S)\ge\delta\mid A_S\}
 \le-\eta_\delta.
\]
If $\mathcal H_S=\{g_*\}$, the posterior on full Boolean functions
converges in total variation to $\delta_{g_*}$. The fixed Gaussian
$\mathcal N(0,K/t_N)$, conditioned on the same $A_S$, has the same
limit. Thus their Boolean posteriors agree asymptotically throughout
$0<\gamma<1$ whenever the compatible energy minimum is unique.
\end{proposition}

\begin{proof}
Couple the output to $Y_N=K^{1/2}\xi/\sqrt{t_N}$ by writing
$Z_N=Q_N^{1/2}\xi/\sqrt{t_N}$, with $\xi$ a standard independent
$M$-dimensional Gaussian. Near the positive definite $K$, the matrix
square root is Lipschitz. For fixed $\epsilon,\delta>0$, kernel
concentration and the Gaussian tail therefore bound
\[
 \Pr\{\|Z_N-Y_N\|_2>\delta\}
 \le C_\epsilon e^{-c_\epsilon N}
 +\Pr\{\|\xi\|_2>\delta\sqrt{t_N}/(C\epsilon)\}.
\]
After division by $t_N$ in the logarithm, let $N\to\infty$ and then
$\epsilon\downarrow0$. Since $t_N=o(N)$, this proves exponential
equivalence at speed $t_N$. Hence $Z_N$ has the Gaussian large-deviation
rate $I_K$.

The minimum rate on the strict orthant equals that on its closure:
multiply its minimizing vector by $1+\varepsilon$ and let
$\varepsilon\downarrow0$. The full-margin denominator consequently
has logarithmic rate $E_S$. Coercivity of $I_K$ and the finiteness of
$\mathcal G_S$ give a strictly positive rate gap on the closed feasible
set at distance at least $\delta$ from $V_S$. The large-deviation upper
bound and the matching denominator lower bound give the displayed
conditional estimate. If the minimizing labeling is unique, the same
argument gives an exponential gap to every other compatible labeling.
Apply the identical argument to $Y_N$.
\end{proof}

If several labelings minimize the energy, this proposition fixes their
possible support but leaves their relative posterior weights undetermined.
Those weights can decide whether a correct label has probability greater
than $2/3$. The next result controls them on a smaller exponent range.

\subsection{Relative probabilities and total variation}

\begin{lemma}[Covariance sensitivity of an orthant probability]
\label{lem:cmp-log-lipschitz}
Let $p_g(Q,t)=\Pr\{\mathcal N(0,Q/t)\in O_g\}$. In a sufficiently
small fixed convex neighborhood of $K$, uniformly over $g$ and $t\ge1$,
\[
 |\log p_g(Q,t)-\log p_g(Q',t)|\le Ct\|Q-Q'\|.
\]
\end{lemma}

\begin{proof}
Uniformly over this covariance neighborhood, integrating over a fixed
box inside each orthant gives $p_g(Q,t)\ge e^{-C_1t}$, with constants
chosen for all finitely many $g$. An unconditional Gaussian tail bound
gives
\[
 \Pr\{\|Z\|_2^2>u\}\le C e^{-ctu}
\]
for $u$ above a fixed constant. Divide by the preceding lower bound
and integrate the resulting conditional tail, starting at a sufficiently
large fixed $u$. It follows that
$\mathbb E[\|Z\|_2^2\mid Z\in O_g]\le C$ uniformly for $t\ge1$.
Differentiating the Gaussian orthant integral along the segment from
$Q$ to $Q'$ now gives a determinant derivative bounded by
$C\|Q-Q'\|$ and a conditional quadratic-moment derivative bounded by
$Ct\|Q-Q'\|$. Integration along the segment proves the claim.
\end{proof}

\begin{theorem}[Quantitative Gaussian comparison]
\label{thm:cmp-relative}
For $0<\gamma<1$, put $t=t_N$. There are positive constants $C,c$
such that
\[
 \max_g\left|
 \frac{\Pr\{Z_N\in O_g\}}{p_g(\widehat K_N,t)}-1\right|
 \le C\frac{t}{\sqrt N}\exp(Ct^2/N)+C\exp(-cN+Ct)
\]
for all sufficiently large $N$. In particular, if $0<\gamma<1/2$,
the Boolean posterior of $Z_N$ conditional on $A_S$ is $o(1)$ in total
variation from that of
$\mathcal N(0,\widehat K_N/t_N)$ conditional on $A_S$.
If $0<\gamma<1/3$, the same statement holds with $K$ in place of
$\widehat K_N$. These assertions include tied minimum-energy labelings.
\end{theorem}

\begin{proof}
Put $D=\|Q_N-\widehat K_N\|$ and restrict first to $D\le u_0$,
with $u_0$ chosen small enough for
Lemma~\ref{lem:cmp-log-lipschitz}. On that event,
\[
 \left|\frac{p_g(Q_N,t)}{p_g(\widehat K_N,t)}-1\right|
 \le e^{CtD}-1.
\]
Integration of the sub-Gaussian tail
\eqref{eq:cmp-finite-concentration} yields
\[
 \mathbb E[(e^{CtD}-1)\mathbf 1_{D\le u_0}]
 \le C\frac{t}{\sqrt N}\exp(Ct^2/N).
\]
For example, integrate
$Ct\int_0^{u_0}e^{Ctu}\Pr(D>u)\,du$ and complete the square;
enlarging the integral to $[0,\infty)$ gives the displayed bound.
On $D>u_0$, the numerator is at most one, its probability is
$Ce^{-cN}$, and the reference orthant probability is at least
$e^{-Ct}$. This proves the relative-probability estimate, including
the negligible subtraction needed in its lower bound.

When $t=o(\sqrt N)$, the estimate tends to zero uniformly over all
full sign patterns. Summing over $\mathcal G_S$ and dividing by the
corresponding sums shows that conditional Boolean probabilities agree
in total variation, uniformly over $S$. Finally,
Lemmas~\ref{lem:cmp-concentration} and
\ref{lem:cmp-log-lipschitz} show that replacing $\widehat K_N$ by
$K$ changes each orthant probability by a factor
$\exp[O(t/b_N)]$. This factor tends to one when
$t/b_N=N^{(3\gamma-1)/2}\to0$, proving the last assertion.
\end{proof}

The two comparisons address distinct errors. Empirical fluctuations of
the hidden kernel lead to the sufficient condition $\gamma<1/2$.
Replacing the deterministic finite-gain kernel by its saturated limit
introduces a separate bias, for which $\gamma<1/3$ is a sufficient
condition. Conditional kernel convergence and the leading Gaussian
large-deviation rate remain valid over the entire range $\gamma<1$.

\section{A Gaussian rate function need not determine the conditioned posterior}
\label{sec:non-gaussian-tie}

Moderate deviations do justify a Gaussian rate function throughout
$0<\gamma<1$. They do not, however, justify replacing the actual
posterior by a conditioned Gaussian in total variation throughout this
interval. The distinction already occurs with three input bits and one
hidden layer. It affects the criterion that the correct label must
have posterior probability greater than $2/3$.

We first give the estimate that makes the counterexample rigorous.
All vectors in the following lemma have a fixed dimension $m$.

\begin{lemma}[Annealed Gaussian corner probabilities]
\label{lem:annealed-corner}
For every $N$, let $H_{N,1},\ldots,H_{N,N}$ be independent copies of a
random vector $H_N\in[-1,1]^m$. Suppose
\[
 K_N=\mathbb E H_NH_N^T\longrightarrow K>0,
 \qquad H_N\ \Longrightarrow\ H.
\]
Set $Q_N=N^{-1}\sum_i H_{N,i}H_{N,i}^T$, and, conditionally on the
features, let $Z_N\sim\mathcal N(0,t_N^{-1}Q_N)$, where
$t_N\to\infty$ and $t_N/N\to0$.
Fix $a>0$ and $g\in\{-1,1\}^m$. Assume
\[
 g_i(K^{-1}g)_i>0\quad\hbox{for every }i.
\]
Writing $v=K^{-1}g$, define
\[
 E_{N,g}=\frac{a^2}{2}g^TK_N^{-1}g,
 \qquad
 C_g=(2\pi)^{-m/2}(\det K)^{-1/2}
             \prod_{i=1}^m\bigl[a g_i v_i\bigr]^{-1}.
\]
Then
\begin{align}
 \log\Pr\{g_iZ_{N,i}>a\ \forall i\}
 ={}&-t_NE_{N,g}-\frac m2\log t_N+\log C_g\notag\\
 &+\frac{a^4t_N^2}{8N}
       \operatorname{Var}\bigl((v^TH)^2\bigr)
       +o\bigl(1+t_N^2/N\bigr).
 \label{eq:annealed-corner}
\end{align}
In particular, the leading correction to the deterministic Gaussian
corner probability can grow without bound even though $t_N=o(N)$.
\end{lemma}

\begin{proof}
Write $D_g=\operatorname{diag}(g)$ and
$E_g(Q)=a^2g^TQ^{-1}g/2$. In a fixed neighborhood of $K$, the
matrices $Q$ are positive definite and all entries of
$\beta_g(Q)=aD_gQ^{-1}g$ are bounded below by a positive constant.
The change of variables $z=ag+D_gu/t_N$ in the Gaussian density gives
\begin{align*}
 \Pr_{Q}\{g_iZ_i>a\ \forall i\}
 ={}&(2\pi)^{-m/2}(\det Q)^{-1/2}t_N^{-m/2}e^{-t_NE_g(Q)}\\
 &\quad\times\int_{[0,\infty)^m}
  \exp\left\{-\beta_g(Q)^Tu
       -\frac{u^TD_gQ^{-1}D_gu}{2t_N}\right\}\,du.
\end{align*}
The integral equals $\prod_i\beta_g(Q)_i^{-1}[1+O(t_N^{-1})]$
uniformly in this neighborhood: use
$0\le1-e^{-r}\le r$ and the finite second moments of the product
exponential density. Denote the resulting smooth positive prefactor
by $C_g(Q)$.

Put $v_N=K_N^{-1}g$ and $\Delta_N=Q_N-K_N$. Matrix inversion has the
uniform expansion
\[
 -t_NE_g(Q_N)
 =-t_NE_{N,g}+t_N\ell_N(\Delta_N)
                  +O(t_N\|\Delta_N\|^2),
 \qquad
 \ell_N(\Delta)=\frac{a^2}{2}v_N^T\Delta v_N.
\]
The linear statistic is the average of the bounded centered variables
\[
 Y_{N,i}=\frac{a^2}{2}
   \left[(v_N^TH_{N,i})^2-
                   \mathbb E(v_N^TH_N)^2\right].
\]
Consequently, with $\theta_N=t_N/N$,
\begin{align*}
 \log\mathbb E e^{t_N\ell_N(\Delta_N)}
 &=N\log\mathbb E e^{\theta_NY_{N,1}}\\
 &=\frac{t_N^2}{2N}\operatorname{Var}(Y_{N,1})
                              +O(t_N^3/N^2).
\end{align*}
All bounds here are uniform, since the $Y_{N,i}$ are uniformly bounded.

For completeness, the quadratic remainder is negligible at the
claimed precision, even when $t_N^2/N$ diverges. Tilt the independent
feature rows by the density proportional to
$\exp(\theta_NY_{N,i})$. Under this product law their mean covariance
differs from $K_N$ by $O(\theta_N)$, and Hoeffding's inequality still
gives, for a fixed number of matrix entries,
\[
 \Pr_{\theta_N}\{\|\Delta_N-
       \mathbb E_{\theta_N}\Delta_N\|>r\}
       \le C e^{-cNr^2}.
\]
Integrating this bound, and using $t_N/N\to0$, yields
\[
 \log\mathbb E_{\theta_N}
   \exp\{C t_N\|\Delta_N\|^2+C\|\Delta_N\|\}
 =O\left(\frac{t_N^3}{N^2}+\frac{t_N}{N}
                                      +N^{-1/2}\right).
\]
A matching lower estimate for the negative of the expression follows
from Jensen's inequality. These bounds control both the inversion
remainder and the smooth prefactor $C_g(Q_N)/C_g(K_N)$.

The complement of the fixed covariance neighborhood has probability
$O(e^{-cN})$, by entrywise Hoeffding. Its contribution to the original
orthant probability is at most this quantity. Inside the neighborhood
the calculation gives a probability $\exp[-O(t_N)-O(\log t_N)]$;
hence the omitted contribution is negligible. The same truncation
under the tilted law is harmless, since the tilt is $o(N)$ and the
features are bounded. We have thus obtained the expansion with
$C_g(K_N)$, $v_N$, and $H_N$, and an error
\[
 O\left(t_N^{-1}+N^{-1/2}+t_N/N+t_N^3/N^2\right)
                  =o(1+t_N^2/N).
\]
Finally $C_g(K_N)\to C_g$, and boundedness and weak convergence give
$\operatorname{Var}((v_N^TH_N)^2)\to
\operatorname{Var}((v^TH)^2)$. This proves the displayed claim.
\end{proof}

\begin{theorem}[An actual-network counterexample to Gaussian posterior equivalence]
\label{thm:gaussian-posterior-counterexample}
Use the original architecture with $d=3$, one hidden layer ($L=2$),
activation $\tanh$, and any fixed margin $a>0$. Order the cube inputs
lexicographically, starting with $x_0=(-1,-1,-1)$. Define
\[
 g=(-1,1,1,-1,-1,-1,-1,1),\qquad
 h=(1,1,1,-1,-1,-1,-1,1).
\]
Observe the seven common labels at every input other than $x_0$.
Let $p_N$ be the actual posterior probability of $g$, after the
all-input-margin restriction and these observations. Let
$\widetilde p_N$ be the corresponding probability for the deterministic
Gaussian output vector $\mathcal N(0,N^{-\gamma}K_N)$, where $K_N$
is the exact mean feature covariance at width $N$. Then
\[
 \widetilde p_N\longrightarrow \frac{189}{314}
                  \qquad (0<\gamma<1),
\]
whereas
\[
 p_N\longrightarrow
 \begin{cases}
  189/314,&0<\gamma<1/2,\\[2pt]
  \displaystyle\left[1+\frac{125}{189}
              e^{-8586a^4/625}\right]^{-1},&\gamma=1/2,\\[6pt]
  1,&1/2<\gamma<1.
 \end{cases}
\]
The same Gaussian limit holds if $K_N$ is replaced by the limiting
arcsine kernel $K$. Thus, for $1/2\le\gamma<1$, the actual Boolean
posterior need not converge in total variation to either of these
conditioned Gaussian posteriors.

For the original margin $a=1/2$ and target $g$, the Gaussian learner's
correct-label probability at the unseen input tends to
$189/314\simeq0.602$, whereas the actual learner's probability tends
to approximately $0.781$ at $\gamma=1/2$ and to one for
$\gamma>1/2$. The proposed $2/3$ confidence criterion therefore
separates their predictions.
\end{theorem}

\begin{proof}
Let $A$ be the $8\times4$ matrix with rows $(1,x)$, in the specified
order; then $A^TA=8I_4$. A feature row has distribution
\[
 H_N=\tanh\left(\frac{N^{(1-\gamma)/2}}{\sqrt3}AW\right),
                 \qquad W\sim\mathcal N(0,I_4).
\]
For $\gamma<1$ it converges almost surely to $H=\operatorname{sign}(AW)$.
Gaussian symmetry and oddness show, exactly at every finite width,
\[
 K_N=(q_N-2c_N)I_8+\frac{c_N}{2}AA^T,
\]
where $q_N$ is the diagonal feature covariance and $c_N$ is the
covariance between cube neighbors. The other off-diagonal entries are
zero at Hamming distance two and $-c_N$ at distance three. Moreover
$q_N\to1$ and $c_N\to1/3$, so
\[
 K=\frac13I_8+\frac16AA^T,\qquad
 K^{-1}=3I_8-\frac3{10}AA^T.
\]
The two functions satisfy
\[
 A^Tg=(-2,-2,2,2)^T,\qquad A^Th=(0,-4,0,0)^T.
\]
Their squared lengths are both eight and the displayed projection
vectors both have squared length sixteen. Therefore
$g^TK_N^{-1}g=h^TK_N^{-1}h$ exactly for all sufficiently large $N$;
the finite-gain correction cannot break the energy tie.

Direct multiplication gives
\begin{align*}
 v:=K^{-1}g&=\frac35(-3,5,5,-7,-1,-3,-3,5),\\
 w:=K^{-1}h&=\frac35(3,3,3,-7,-3,-3,-3,7).
\end{align*}
All signed gradient coordinates are strictly positive, and
$g^Tv=h^Tw=96/5$. Both orthant minimizers are therefore their corners
$ag$ and $ah$. The Gaussian corner prefactors have ratio
\[
 \frac{C_g}{C_h}
 =\frac{\prod_i h_i w_i}{\prod_i g_i v_i}
 =\frac{3^6\,7^2}{3^3\,5^3\,7}=\frac{189}{125}.
\]
The uniform Gaussian corner estimate in
Lemma~\ref{lem:annealed-corner}, before averaging over features, proves
$\widetilde p_N\to189/314$.

It remains to compute a fourth moment of the saturated feature row.
Here this calculation is exact and finite. Signed permutations of the
four coordinates of $W$ give three orbits of possible threshold
vectors $H$, of sizes $8,64,32$. To check completeness, first sort the
absolute coordinates as $r_1>r_2>r_3>r_4>0$. Only the signs of
$r_1-r_2-r_3-r_4$ and $r_1-r_2-r_3+r_4$ are undetermined; their
three feasible sign combinations give precisely these three orbits.
The exceptional equalities have probability zero.

Let $P=AA^T/8$. The following table records finite averages within
each orbit. Representatives are coefficient vectors $W$, not feature
vectors. The column $n_+$ counts orbit elements with coordinates
$H_1,H_2,H_4,H_7$ all positive, using zero-based cube indices.
\[
\begin{array}{c|r|r|r|r|r}
 \text{representative }W&\text{size}&\|PH\|^2&n_+
 &\text{mean }(v^TH)^4&\text{mean }(w^TH)^4\\ \hline
 (10,3,2,1)&8&8&1&1296/625&5184/625\\
 (5,4,2,3/2)&64&6&4&964224/625&1174176/625\\
 (4,7/2,3,1)&32&6&0&1174176/625&754272/625
\end{array}
\]
Each entry can be checked by applying the $2^4 4!$ signed
permutations to the representative, taking $\operatorname{sign}(AW)$,
removing duplicates, and summing the indicated integer fourth powers.
Gaussian invariance makes the distribution uniform within each orbit.
To determine the three orbit probabilities $p_1,p_2,p_3$, note that
\[
 8p_1+6p_2+6p_3
       =\mathbb E\|PH\|^2=\operatorname{tr}(PK)=20/3.
\]
Thus $p_1=1/3$. The four rows of $A$ with indices $1,2,4,7$ are
orthogonal, so their Gaussian projections have independent signs.
Using the $n_+$ column,
\[
 \frac1{16}=\Pr\{H_1=H_2=H_4=H_7=1\}
            =\frac{p_1}{8}+\frac{p_2}{16}.
\]
It follows that $p_2=p_3=1/3$. The table consequently gives
\[
 \mathbb E(v^TH)^4=\frac{713232}{625},\qquad
 \mathbb E(w^TH)^4=\frac{644544}{625}.
\]
The second moments are equal, both being $96/5$. Hence the difference
of the two variances of squared projections is $68688/625$.

Apply Lemma~\ref{lem:annealed-corner} with $t_N=N^\gamma$.
Since the two energies agree exactly, their prior log odds satisfy
\[
 \log\frac{\Pr(A_g)}{\Pr(A_h)}
 =\log\frac{189}{125}
    +\frac{8586a^4}{625}N^{2\gamma-1}
    +o(1+N^{2\gamma-1}).
\]
The seven observed labels leave exactly these two Boolean functions
compatible, so these are also their posterior odds. The three stated
limits follow directly.
\end{proof}

\begin{remark}[What this counterexample does and does not show]
The leading kernel margin energy remains correct at every
$0<\gamma<1$: both functions in the example have the same energy, and
the actual posterior selects one of these tied minimizers. The
counterexample disproves full conditional-Gaussian equivalence, even
with the exact finite-width mean covariance. It does not by itself
establish a circuit prior or a general sample-complexity advantage.
It shows that the leading kernel energy alone need not determine the
posterior's confidence, and that fourth-order feature statistics can
remain visible after rare-event conditioning.
\end{remark}


\section{Consequences for the circuit-prior research question}

The strongest general statements proved here are as follows.
For $\gamma\le0$, the limiting posterior is the Gaussian sign
posterior, with a global margin at $\gamma=0$ and no margin at
$\gamma<0$. For $0<\gamma<1/2$, the deterministic finite-gain
Gaussian approximates every full Boolean posterior, including tied
energy minima. The sufficient range for replacing that kernel by
its saturated limit is $0<\gamma<1/3$; optimality of the latter
threshold is not asserted. For every $0<\gamma<1$, the saturated
Gaussian rate and concentration on compatible energy minimizers
remain valid. A unique compatible minimum gives the same limiting
Boolean posterior throughout this entire interval. With ties,
Theorem~\ref{thm:gaussian-posterior-counterexample} disproves a
universal Gaussian-posterior replacement for every
$1/2\le\gamma<1$.

These comparisons concern probabilities themselves. A strict
$2/3$ confidence indicator is discontinuous at probability exactly
$2/3$; total-variation convergence alone does not settle that
boundary case. The counterexample stays strictly away from the
boundary on both sides and therefore gives an unambiguous
separation under the requested convention.

A companion paper, \emph{Small circuits can have exponentially
large prior complexity below the mean-field scaling}, proves a
separate conclusion throughout $\gamma<1$: full parity has a
polynomial-size, polynomial-weight $\tanh$ circuit but exponential
prior complexity on the natural scale of the regime. That prior
separation does not imply exponential labeled-sample complexity.
Conversely, the present non-Gaussian correction does not restore a
polynomial circuit prior. It resolves a different issue: leading
kernel energies need not determine how the Bayesian learner
weights equally costly explanations.

The remaining statistical question is whether the actual
all-input-margin-conditioned posterior needs exponentially many
examples to predict a $k$-bit XOR correctly with probability
strictly greater than $2/3$ on most inputs. Neither the prior-mass
separation nor the usual posterior-mean regression argument proves
that statement. The counterexample also means that a proof covering
the whole subcritical interval cannot simply replace every posterior
by a Gaussian posterior and ignore its tied minima.

\paragraph{Verification.}
The accompanying script, \path{verify-counterexample.py}, checks the finite orbit table,
projection identities, strict corner gradients, energy equality,
Gaussian prefactor ratio, and fourth-moment coefficient using exact
rational arithmetic. The asymptotic estimates and the completeness
of the orbit classification are proved in the text.

\end{document}
